\documentclass[10pt,a4paper]{amsart}
\usepackage{amsmath,amssymb,amsthm,graphicx}
\usepackage{verbatim}
\newenvironment{reference}{\comment}{\endcomment}
%\newenvironment{reference}{}{}
\newenvironment{slogan}{\comment}{\endcomment}
\newenvironment{history}{\comment}{\endcomment}

% For commutative diagrams we use Xy-pic
\usepackage[all]{xy}

% We use 2cell for 2-commutative diagrams.
\xyoption{2cell}
% Theorem environments.
%
\theoremstyle{plain}
\newtheorem{theorem}[subsection]{Theorem}
\newtheorem{proposition}[subsection]{Proposition}
\newtheorem{lemma}[subsection]{Lemma}

\theoremstyle{definition}
\newtheorem{definition}[subsection]{Definition}
\newtheorem{example}[subsection]{Example}
\newtheorem{exercise}[subsection]{Exercise}
\newtheorem{situation}[subsection]{Situation}

\theoremstyle{remark}
\newtheorem{remark}[subsection]{Remark}
\newtheorem{remarks}[subsection]{Remarks}

\numberwithin{equation}{subsection}

% Macros
%
\def\lim{\mathop{\mathrm{lim}}\nolimits}
\def\colim{\mathop{\mathrm{colim}}\nolimits}
\def\Spec{\mathop{\mathrm{Spec}}}
\def\Hom{\mathop{\mathrm{Hom}}\nolimits}
\def\Ext{\mathop{\mathrm{Ext}}\nolimits}
\def\SheafHom{\mathop{\mathcal{H}\!\mathit{om}}\nolimits}
\def\SheafExt{\mathop{\mathcal{E}\!\mathit{xt}}\nolimits}
\def\Sch{\mathit{Sch}}
\def\Mor{\mathop{\mathrm{Mor}}\nolimits}
\def\Ob{\mathop{\mathrm{Ob}}\nolimits}
\def\Sh{\mathop{\mathit{Sh}}\nolimits}
\def\NL{\mathop{N\!L}\nolimits}
\def\CH{\mathop{\mathrm{CH}}\nolimits}
\def\proetale{{pro\text{-}\acute{e}tale}}
\def\etale{{\acute{e}tale}}
\def\QCoh{\mathit{QCoh}}
\def\Ker{\mathop{\mathrm{Ker}}}
\def\Im{\mathop{\mathrm{Im}}}
\def\Coker{\mathop{\mathrm{Coker}}}
\def\Coim{\mathop{\mathrm{Coim}}}

% Boxtimes
%
\DeclareMathSymbol{\boxtimes}{\mathbin}{AMSa}{"02}

%
% Macros for moduli stacks/spaces
%
\def\QCohstack{\mathcal{QC}\!\mathit{oh}}
\def\Cohstack{\mathcal{C}\!\mathit{oh}}
\def\Spacesstack{\mathcal{S}\!\mathit{paces}}
\def\Quotfunctor{\mathrm{Quot}}
\def\Hilbfunctor{\mathrm{Hilb}}
\def\Curvesstack{\mathcal{C}\!\mathit{urves}}
\def\Polarizedstack{\mathcal{P}\!\mathit{olarized}}
\def\Complexesstack{\mathcal{C}\!\mathit{omplexes}}
% \Pic is the operator that assigns to X its picard group, usage \Pic(X)
% \Picardstack_{X/B} denotes the Picard stack of X over B
% \Picardfunctor_{X/B} denotes the Picard functor of X over B
\def\Pic{\mathop{\mathrm{Pic}}\nolimits}
\def\Picardstack{\mathcal{P}\!\mathit{ic}}
\def\Picardfunctor{\mathrm{Pic}}
\def\Deformationcategory{\mathcal{D}\!\mathit{ef}}
\usepackage{cleveref}
\crefname{originallemma}{original lemma}{original lemmas}
\Crefname{originallemma}{Original lemma}{Original lemmas}
\makeatletter
\newlabel{lemma-cocontinuous-morphism-topoi}{{00XO}{source}{}{sourcepointer-1}{}}
\newlabel{lemma-cocontinuous-morphism-topoi@cref}{{[originallemma][0][]00XO}{[1][0][]source}{}{}{}}
\makeatother
\crefname{originallemma}{original lemma}{original lemmas}
\Crefname{originallemma}{Original lemma}{Original lemmas}
\makeatletter
\newlabel{lemma-when-shriek}{{00XR}{source}{}{sourcepointer-2}{}}
\newlabel{lemma-when-shriek@cref}{{[originallemma][0][]00XR}{[1][0][]source}{}{}{}}
\makeatother
\crefname{originallemma}{original lemma}{original lemmas}
\Crefname{originallemma}{Original lemma}{Original lemmas}
\makeatletter
\newlabel{lemma-adjoint-functors}{{09VQ}{source}{}{sourcepointer-3}{}}
\newlabel{lemma-adjoint-functors@cref}{{[originallemma][0][]09VQ}{[1][0][]source}{}{}{}}
\makeatother
\crefname{originallemma}{original lemma}{original lemmas}
\Crefname{originallemma}{Original lemma}{Original lemmas}
\makeatletter
\newlabel{categories-lemma-adjoint-exact}{{0038}{source}{}{sourcepointer-4}{}}
\newlabel{categories-lemma-adjoint-exact@cref}{{[originallemma][0][]0038}{[1][0][]source}{}{}{}}
\makeatother
\crefname{originallemma}{original lemma}{original lemmas}
\Crefname{originallemma}{Original lemma}{Original lemmas}
\makeatletter
\newlabel{lemma-limit-sheaf}{{00W2}{source}{}{sourcepointer-5}{}}
\newlabel{lemma-limit-sheaf@cref}{{[originallemma][0][]00W2}{[1][0][]source}{}{}{}}
\makeatother
\crefname{originallemma}{original lemma}{original lemmas}
\Crefname{originallemma}{Original lemma}{Original lemmas}
\makeatletter
\newlabel{lemma-adjoints-u}{{00VE}{source}{}{sourcepointer-6}{}}
\newlabel{lemma-adjoints-u@cref}{{[originallemma][0][]00VE}{[1][0][]source}{}{}{}}
\makeatother
\crefname{originallemma}{original lemma}{original lemmas}
\Crefname{originallemma}{Original lemma}{Original lemmas}
\makeatletter
\newlabel{lemma-sheafification-exact}{{00WJ}{source}{}{sourcepointer-7}{}}
\newlabel{lemma-sheafification-exact@cref}{{[originallemma][0][]00WJ}{[1][0][]source}{}{}{}}
\makeatother
\crefname{originallemma}{original lemma}{original lemmas}
\Crefname{originallemma}{Original lemma}{Original lemmas}
\makeatletter
\newlabel{lemma-pushforward-sheaf}{{00WW}{source}{}{sourcepointer-8}{}}
\newlabel{lemma-pushforward-sheaf@cref}{{[originallemma][0][]00WW}{[1][0][]source}{}{}{}}
\makeatother
\crefname{originallemma}{original lemma}{original lemmas}
\Crefname{originallemma}{Original lemma}{Original lemmas}
\makeatletter
\newlabel{lemma-adjoint-sheaves}{{00WX}{source}{}{sourcepointer-9}{}}
\newlabel{lemma-adjoint-sheaves@cref}{{[originallemma][0][]00WX}{[1][0][]source}{}{}{}}
\makeatother
\crefname{originallemma}{original lemma}{original lemmas}
\Crefname{originallemma}{Original lemma}{Original lemmas}
\makeatletter
\newlabel{lemma-have-functor-other-way}{{00XX}{source}{}{sourcepointer-10}{}}
\newlabel{lemma-have-functor-other-way@cref}{{[originallemma][0][]00XX}{[1][0][]source}{}{}{}}
\makeatother
\title{Consequences of the existence of a left adjoint for a cocontinuous functor}
\author{The Stacks Project Authors}
\date{stacks-project:sites.tex:local-snapshot-20261012; GNU FDL 1.2 or later}
\begin{document}
\maketitle
\noindent\textit{Editorial note.} All delivered mathematical text is the author's own, copied from the pinned Stacks Project chapter file sites-cohomology.tex; no statement or proof step was rewritten, reordered or removed.\par\medskip
\noindent\textit{Editorial note.} The proof is the author's complete argument; every step the author writes out is inside the counted span and no closing sentence is a placeholder.\par\medskip
\noindent\textit{Editorial note.} Every labelled reference occurring inside the counted spans is carried in assembly.external\_references as a non-load-bearing navigation pointer (same-chapter ones at their definition site, cross-chapter ones at the invocation inside the counted span), so no load-bearing mathematics is replaced by a stub. No bibliography citation occurs inside the counted spans.\par\medskip
\noindent\textit{Editorial note.} The checked chapter list for same-name collisions is the set of stacks chapters staged in this repository: chow, cohomology, divisors, etale-cohomology, flat, more-algebra, more-morphisms, morphisms, spaces-more-morphisms, varieties together with this card's chapter sites-cohomology.\par\medskip
\noindent\textit{Editorial note.} The preamble is the chapter's own retained minimal preamble (sources/stacks-sites-cohomology/retained-preamble.tex), a verbatim excerpt of the native preamble: the Xy-pic setup and the native theorem declarations with the mathematical macros.\par\medskip
\noindent\textit{Allowed standard prerequisites (editorial).}\begin{itemize}
\item The functors on presheaves attached to an adjoint pair. Same chapter, Lemma lemma-adjoint-functors and Lemma lemma-adjoints-u. Identify the left Kan extension with the direct image along the left adjoint and show that it is a right adjoint.
\item Limits of sheaves and exactness of sheafification. Same chapter, Lemma lemma-limit-sheaf and Lemma lemma-sheafification-exact. Transfer the commutation with finite limits from presheaves to sheaves.
\item The pushforward along a continuous functor and its adjoint. Same chapter, Lemma lemma-pushforward-sheaf and Lemma lemma-adjoint-sheaves. Supply the left adjoint in the continuous case.
\end{itemize}\medskip
\ifdefined\hypertarget\hypertarget{sourcepointer-1}{}\fi
\noindent\textit{Source reference.} Original lemma 00XO: sources/stacks-sites/pinned-original-source.tex, source lines 5205--5205. Navigation pointer invoked at proof line 5205; the counted span contains the pointer only and the cited result is not reproduced inside it, so the pointer is registered as a non-load-bearing navigation stub in the established form for Stacks material. No bibliography citation is involved.\par\medskip
\ifdefined\hypertarget\hypertarget{sourcepointer-2}{}\fi
\noindent\textit{Source reference.} Original lemma 00XR: sources/stacks-sites/pinned-original-source.tex, source lines 5206--5206. Navigation pointer invoked at proof line 5206; the counted span contains the pointer only and the cited result is not reproduced inside it, so the pointer is registered as a non-load-bearing navigation stub in the established form for Stacks material. No bibliography citation is involved.\par\medskip
\ifdefined\hypertarget\hypertarget{sourcepointer-3}{}\fi
\noindent\textit{Source reference.} Original lemma 09VQ: sources/stacks-sites/pinned-original-source.tex, source lines 5225--5225. Navigation pointer invoked at proof line 5225; the counted span contains the pointer only and the cited result is not reproduced inside it, so the pointer is registered as a non-load-bearing navigation stub in the established form for Stacks material. No bibliography citation is involved.\par\medskip
\ifdefined\hypertarget\hypertarget{sourcepointer-4}{}\fi
\noindent\textit{Source reference.} Original lemma 0038: sources/stacks-sites/pinned-original-source.tex, source lines 5229--5229. Navigation pointer invoked at proof line 5229; the counted span contains the pointer only and the cited result is not reproduced inside it, so the pointer is registered as a non-load-bearing navigation stub in the established form for Stacks material. No bibliography citation is involved.\par\medskip
\ifdefined\hypertarget\hypertarget{sourcepointer-5}{}\fi
\noindent\textit{Source reference.} Original lemma 00W2: sources/stacks-sites/pinned-original-source.tex, source lines 5232--5232. Navigation pointer invoked at proof line 5232; the counted span contains the pointer only and the cited result is not reproduced inside it, so the pointer is registered as a non-load-bearing navigation stub in the established form for Stacks material. No bibliography citation is involved.\par\medskip
\ifdefined\hypertarget\hypertarget{sourcepointer-6}{}\fi
\noindent\textit{Source reference.} Original lemma 00VE: sources/stacks-sites/pinned-original-source.tex, source lines 5233--5233. Navigation pointer invoked at proof line 5233; the counted span contains the pointer only and the cited result is not reproduced inside it, so the pointer is registered as a non-load-bearing navigation stub in the established form for Stacks material. No bibliography citation is involved.\par\medskip
\ifdefined\hypertarget\hypertarget{sourcepointer-7}{}\fi
\noindent\textit{Source reference.} Original lemma 00WJ: sources/stacks-sites/pinned-original-source.tex, source lines 5236--5236. Navigation pointer invoked at proof line 5236; the counted span contains the pointer only and the cited result is not reproduced inside it, so the pointer is registered as a non-load-bearing navigation stub in the established form for Stacks material. No bibliography citation is involved.\par\medskip
\ifdefined\hypertarget\hypertarget{sourcepointer-8}{}\fi
\noindent\textit{Source reference.} Original lemma 00WW: sources/stacks-sites/pinned-original-source.tex, source lines 5242--5242. Navigation pointer invoked at proof line 5242; the counted span contains the pointer only and the cited result is not reproduced inside it, so the pointer is registered as a non-load-bearing navigation stub in the established form for Stacks material. No bibliography citation is involved.\par\medskip
\ifdefined\hypertarget\hypertarget{sourcepointer-9}{}\fi
\noindent\textit{Source reference.} Original lemma 00WX: sources/stacks-sites/pinned-original-source.tex, source lines 5242--5242. Navigation pointer invoked at proof line 5242; the counted span contains the pointer only and the cited result is not reproduced inside it, so the pointer is registered as a non-load-bearing navigation stub in the established form for Stacks material. No bibliography citation is involved.\par\medskip
\ifdefined\hypertarget\hypertarget{sourcepointer-10}{}\fi
\noindent\textit{Source reference.} Original lemma 00XX: sources/stacks-sites/pinned-original-source.tex, source lines 5143--5143. Navigation pointer invoked at proof line 5143; the counted span contains the pointer only and the cited result is not reproduced inside it, so the pointer is registered as a non-load-bearing navigation stub in the established form for Stacks material. No bibliography citation is involved.\par\medskip
\begin{lemma}
\label{lemma-have-left-adjoint}
Let $\mathcal{C}$ and $\mathcal{D}$ be sites. Let
$g : \Sh(\mathcal{C}) \to \Sh(\mathcal{D})$ be
the morphism of topoi associated to a continuous and cocontinuous functor
$u : \mathcal{C} \to \mathcal{D}$, see
Lemmas \ref{lemma-cocontinuous-morphism-topoi} and
\ref{lemma-when-shriek}.
\begin{enumerate}
\item If $w : \mathcal{D} \to \mathcal{C}$ is a left adjoint to $u$, then
\begin{enumerate}
\item $g_!\mathcal{F}$ is the sheaf associated to the presheaf
$w^p\mathcal{F}$, and
\item $g_!$ is exact.
\end{enumerate}
\item if $w$ is a continuous left adjoint, then $g_!$
has a left adjoint.
\item If $w$ is a cocontinuous left adjoint, then $g_! = h^{-1}$ and
$g^{-1} = h_*$ where $h : \Sh(\mathcal{D}) \to \Sh(\mathcal{C})$ is
the morphism of topoi associated to $w$.
\end{enumerate}
\end{lemma}

\begin{proof}
Recall that $g_!\mathcal{F}$ is the sheafification of $u_p\mathcal{F}$.
Hence (1)(a) follows from the fact that $u_p = w^p$ by
Lemma \ref{lemma-adjoint-functors}.

\medskip\noindent
To see (1)(b) note that $g_!$ commutes with all colimits as $g_!$
is a left adjoint (Categories, Lemma \ref{categories-lemma-adjoint-exact}).
Let $i \mapsto \mathcal{F}_i$ be a finite diagram in $\Sh(\mathcal{C})$.
Then $\lim \mathcal{F}_i$ is computed in the category of presheaves
(Lemma \ref{lemma-limit-sheaf}). Since $w^p$ is a right adjoint
(Lemma \ref{lemma-adjoints-u})
we see that $w^p \lim \mathcal{F}_i = \lim w^p\mathcal{F}_i$. Since
sheafification is exact
(Lemma \ref{lemma-sheafification-exact})
we conclude by (1)(a).

\medskip\noindent
Assume $w$ is continuous. Then $g_! = (w^p\ )^\# = w^s$ but sheafification
isn't necessary and one has the left adjoint $w_s$, see
Lemmas \ref{lemma-pushforward-sheaf} and \ref{lemma-adjoint-sheaves}.

\medskip\noindent
Assume $w$ is cocontinuous. The equality $g_! = h^{-1}$ follows from (1)(a)
and the definitions. The equality $g^{-1} = h_*$ follows from the equality
$g_! = h^{-1}$ and uniqueness of adjoint functor. Alternatively one can deduce
it from Lemma \ref{lemma-have-functor-other-way}.
\end{proof}

\end{document}
